% ---
% titre: Calcul mental
% theme: NO
% chapitre: Nombres naturels et décimaux
% sous_chapitre: PPMC et PGDC
% niveau: [9e]
% type: EVAL
% tags: [c-nombres-premiers, c-decomposition-facteurs-premiers, p-question-TS]
% difficulte: 1
% corrige: true
% ---

\question
Détermine:

\begin{parts} 

\vspace{20pt}\part[6]
le ppmc et le pgdc de 18 et 30 par la méthode de la liste.

\vspace{20pt}{\footnotesize\raggedleft Espace pour ta démarche et tes calculs\par}
	\begin{solutionorgrid}[160pt]
		M$(18)$ = $\{18; 36; 54; 72; 90; 108; 126; \ldots\}$
		
		M$(30)$ = $\{30; 60; 90; 120; \ldots\}$
		
		\vspace{10pt}
		ppmc(18;30) $= 90$
		
		\vspace{10pt}
		\textbf{1pt} pour  chaque liste de multiples
		
		\textbf{1pt} pour le ppmc
		
		\vspace{20pt}
		D$(18)$ = $\{1; 2; 3; 6; 9; 18\}$
		
		D$(30)$ = $\{1; 2; 3; 5; 6; 10; 15; 30\}$
		
		\vspace{10pt}
		pgdc(18;30) $= 6$
		
		\vspace{10pt}
		\textbf{1pt} pour chaque liste de diviseurs
		
		\textbf{1pt} pour le pgdc
	\end{solutionorgrid}

\vspace{20pt}\part[4]
le ppmc et le pgdc de 88 et 132 par la méthode de la décomposition en facteurs premiers.

\vspace{20pt}{\footnotesize\raggedleft Espace pour ta démarche et tes calculs\par}
	\begin{solutionorgrid}[160pt]
		\begin{minipage}[t]{0.12\linewidth}
		\vspace{0pt}
		\begin{tabularx}{\linewidth}{>{\raggedleft\arraybackslash}X|X}
		88&2\\
		44&2\\
		22&2\\
		11&11\\
		1&\\
		\end{tabularx}
		\end{minipage}%
		\hfill
		\begin{minipage}[t]{0.12\linewidth}
		\vspace{0pt}
		\begin{tabularx}{\linewidth}{>{\raggedleft\arraybackslash}X|X}
		132&2\\
		66&2\\
		33&3\\
		11&11\\
		1&\\
		\end{tabularx}
		\end{minipage}%
		\hfill
		\begin{minipage}[t]{0.7\linewidth}
		\vspace{0pt}
		$88 = 2^3 \cdot 11$ \quad / \quad $132 = 2^2 \cdot 3 \cdot 11$
		
		\vspace{10pt}
		ppmc(88;132) $= 2^3\cdot 3 \cdot 11 = 264$
		
		\vspace{10pt}
		pgdc(88;132) $= 2^2\cdot 11 = 44$
				
		\vspace{10pt}
		\textbf{1pt} pour chaque décomposition
		
		\textbf{1pt} pour le ppmc et  \textbf{1pt} pour le pgdc
		\end{minipage}%
	\end{solutionorgrid}
\end{parts}

